\chapter{Introduction}
\label{chap:intro}

The Intensive Care Unit (ICU) represents a demanding clinical environment where
continuous patient observation is paramount, especially for individuals in a
comatose or mobility-impaired state. Current monitoring protocols are strained
by the inherent limits of human vigilance and the shortcomings of existing
automated alarms, which frequently produce high rates of false positives and
contribute to the well-documented phenomenon of ``alarm fatigue.''

This project introduces \textbf{Vital Guardian}, an artificial intelligence-driven
patient monitoring system engineered to bridge this deficiency. Vital Guardian
operates as a vigilant, always-on layer of observation designed to empower
clinical staff by providing timely, contextually aware alerts for critical patient
events. The system intelligently fuses data from two sensory streams --- computer
vision and audio analysis --- and preserves patient privacy by processing all raw
media locally on an edge device.

% =============================================================
\section{Existing Solutions}
\label{sec:existing_solutions}

Patient monitoring in ICUs currently relies on a combination of direct human
supervision and various automated technologies. While these approaches provide a
baseline of safety, each suffers from distinct limitations in scalability,
accuracy, and privacy.

Human attendants provide the most comprehensive care but are subject to fatigue
and staffing shortages. Automated solutions such as pressure mats and wearable
sensors generate high rates of false alarms due to routine patient restlessness
rather than genuine emergencies. Centralised CCTV systems raise significant
privacy concerns by transmitting video data to central servers.

Table \ref{tab:existing_solutions} compares the prevailing monitoring approaches
with the proposed Vital Guardian system.

\begin{table}[H]
    \centering
    \caption{Comparison of Existing Patient Monitoring Solutions}
    \label{tab:existing_solutions}
    \renewcommand{\arraystretch}{1.5}
    \begin{tabular}{|p{3.5cm}|p{5cm}|p{5cm}|}
    \hline
    \textbf{Solution} & \textbf{Overview} & \textbf{Key Limitations} \\ \hline
    Human Attendants &
    Continuous manual observation by nursing staff or dedicated sitters. &
    Not scalable; subject to fatigue; response times limited by
    high patient-to-nurse ratios. \\ \hline
    Bed Pressure Alarms &
    Sensors embedded in the bed that trigger when weight is removed. &
    Extremely high false positive rates; cannot distinguish routine
    repositioning from a genuine fall. \\ \hline
    Centralised CCTV and Cloud AI &
    Video feeds transmitted to a central station or cloud server for analysis. &
    High privacy risk from data transmission; latency issues; requires
    expensive infrastructure upgrades. \\ \hline
    Vital Guardian (Proposed) &
    Multi-modal AI monitoring with all inference performed locally on an edge device. &
    Requires capable edge hardware; Cognitive Core reasoning depends on
    external LLM connectivity. \\ \hline
    \end{tabular}
\end{table}

% =============================================================
\section{Problem Statement}
\label{sec:problem_statement}

Comatose patients in ICUs represent one of the most vulnerable populations in
healthcare, requiring continuous supervision due to their inability to communicate
distress or respond independently to emergencies. The healthcare industry lacks a
comprehensive solution that provides intelligent, automated monitoring while
simultaneously preserving patient privacy and reducing caregiver burden.

This gap in care represents a critical safety issue. Delayed detection of
emergencies such as falls, seizures, or respiratory distress can result in severe
complications or death. Existing automated systems exacerbate the problem through
alarm fatigue, where frequent false alarms desensitise staff to genuine
emergencies. Cloud-based monitoring solutions introduce further unacceptable
privacy risks. Vital Guardian is developed to directly address these issues by
providing a context-aware, privacy-preserving monitoring solution that processes
all raw data on-premise.

% =============================================================
\section{Scope}
\label{sec:scope}

The scope of the Vital Guardian project is to design, develop, and validate an
edge-computing system for intelligent ICU patient monitoring.

\begin{itemize}
    \item \textbf{Inclusions:} The project covers the development of a multi-modal
    monitoring solution integrating computer vision and audio processing. It includes
    real-time patient tracking, fall detection, dangerous movement analysis
    (seizures), audio distress classification (gasps, moans), and multilingual
    keyword spotting. It includes a multi-tier LLM-driven Cognitive Core for
    contextual reasoning and natural language alert generation.

    \item \textbf{Exclusions:} The system does not retain raw video or audio data
    for long-term storage. It does not directly interface with or control
    life-support equipment such as ventilators. It is an assistive tool and does
    not replace the legal requirement for qualified human medical oversight.
\end{itemize}

% =============================================================
\section{Modules}
\label{sec:modules}

The Vital Guardian system is decomposed into three primary modules.

\subsection{Module 1: The Visual Guardian}

This module serves as the eyes of the system, responsible for all vision-based
monitoring tasks using a standard IP or USB camera.

\begin{enumerate}
    \item \textbf{Patient Tracking:} Continuous detection and localisation of the
    patient within the camera frame using YOLOv11n, accelerated via Intel OpenVINO.
    \item \textbf{Fall Detection:} Real-time detection of rapid descents or
    unintended bed exits using a MoViNet-A2 spatiotemporal classifier on a
    16-frame rolling buffer.
    \item \textbf{Seizure Detection:} Detection of prolonged rhythmic movements
    using a dedicated MoViNet-A2 model operating on a 64-frame temporal window
    with stride-2 subsampling.
    \item \textbf{Kinematic Analysis:} Biomechanical verification of detected
    events using MediaPipe Pose landmark tracking.
\end{enumerate}

\subsection{Module 2: The Auditory Watchdog}

This module functions as the ears of the system, providing audio intelligence
to detect non-visual signs of patient distress.

\begin{enumerate}
    \item \textbf{Distress Classification:} Identifies non-verbal distress sounds
    including gasps, moans, and wheezing using YAMNet.
    \item \textbf{Keyword Spotting:} Detects predefined help-related keywords in
    both English and Urdu using Faster-Whisper transcription.
    \item \textbf{Privacy Shield:} Utilises Silero Voice Activity Detection (VAD)
    to identify continuous conversations and automatically pause speech
    transcription, protecting visitor and patient privacy.
\end{enumerate}

\subsection{Module 3: The Cognitive Core}

This module serves as the brain of the system, using Gemini language models to
fuse data from the visual and audio modules into structured clinical alerts.

\begin{enumerate}
    \item \textbf{Tier-2 Verification:} Receives structured JSON payloads from
    both sensory modules and performs binary confirmation or suppression of the
    detected event, reducing false positives.
    \item \textbf{Tier-3 Enrichment:} Generates human-readable, clinically
    actionable alert messages with a severity rating for display on the nursing
    dashboard.
\end{enumerate}

% =============================================================
\section{Work Division}
\label{sec:work_division}

The project follows a structured, role-based implementation approach.

\begin{table}[H]
    \centering
    \caption{Project Work Division and Member Responsibilities}
    \label{tab:work_division}
    \renewcommand{\arraystretch}{1.5}
    \begin{tabular}{|p{4cm}|p{4cm}|p{5.5cm}|}
    \hline
    \textbf{Member Name} & \textbf{Role} & \textbf{Primary Responsibility} \\ \hline
    Eman Hassan &
    Project Lead and System Architect &
    System design, Cognitive Core integration, and overall project planning. \\ \hline
    Muhammad Mahad Khan &
    Vision Module Lead &
    Development of YOLOv11n, MoViNet-A2, and MediaPipe vision pipeline. \\ \hline
    Muhammad Hasan Waqar &
    Audio Module Lead &
    Development of YAMNet distress classification and Keyword Spotting. \\ \hline
    \end{tabular}
\end{table}
